%%%%%%%%%%%%%%%%%%%%%%%%%%%%%%%%%%%%%%%%%%%%%%%%%%%%%%%%%%%%%%%%%%%%%%
% Lớp 11 | Giới hạn. Hàm số liên tục | Giới hạn của dãy số
% Dạng 1D3?1-4: Giới hạn vô cực
% Số câu: 41  (sinh tự động bởi tach_cau.py; từ nay thầy sửa trực tiếp file này)
%%%%%%%%%%%%%%%%%%%%%%%%%%%%%%%%%%%%%%%%%%%%%%%%%%%%%%%%%%%%%%%%%%%%%%

%%%%%%%%%%%%%%%% TRẮC NGHIỆM 4 PHƯƠNG ÁN %%%%%%%%%%%%%%%%

% nguồn: Ok_Toan11_TracNghiem_Lan2
\begin{ex}%[1D3N1-4]%[ma:661eedc5]
	Tính $\lim (-2n^{2\,019}+3n^{2\,018}+4)$. 
	\choice 
		{\True $-\infty$} 
		{$+\infty$} 
		{$-2$} 
		 {$2\,019$}
	 \loigiai{
		Ta có $\lim\left(-2n^{2\,019}+3n^{2018}+4\right)=\lim\left[n^{2\,019}\cdot\left(-2+\dfrac{3}{n}+\dfrac{4}{n^{2\,019}} \right)\right]=-\infty$\\ vì $\heva{&\lim n^{2019}=+\infty  \\&\lim\left( -2+\dfrac{3}{n}+\dfrac{4}{n^{2\,019}}\right)=-2<0.}$
	} 
\end{ex}

% nguồn: Ok_Toan11_TracNghiem_Lan2
\begin{ex}%[1D3N1-4]%[ma:898a9bbb]
	$\lim  {\left(2^{n}-1\right)}$ bằng
	\choice 
		{$-1$} 
		{$1$} 
		{\True  $+\infty$}
		 {$-\infty$}
	\loigiai{ 
		Ta có $\lim\left(2^{n}-1\right)=\lim 2^{n}\left(1-\dfrac{1}{2^{n}}\right)$.\\ 
		Vì $\heva{&\lim  2^{n}=+\infty \\ &\lim\left[1-\left(\dfrac{1}{2}\right)^{n}\right]=1>0}$ nên $\lim\left(2^{n}-1\right)=+\infty$.
	} 
\end{ex}

% nguồn: Ok_Toan11_TracNghiem_Lan2
\begin{ex}%[1D3N1-4]%[ma:d9b2982b]
	Dãy số nào sau đây  có giới hạn bằng $0$?
	\choice 
		{$\left(\dfrac{4}{\mathrm{e}}\right)^{n}$} 
		{\True  $\left(\dfrac{1}{3}\right)^{n}$}
		{$\left(\dfrac{5}{3}\right)^{n}$} 
		{$\left(\dfrac{-5}{3}\right)^{n}$} 
	\loigiai{ 
		Ta có $\lim q^{n}=0$  nếu $|q|<1$.\\
		Mặt khác $\left|\dfrac{4}{\mathrm{e}}\right|>1$; $\left|\dfrac{5}{3}\right|=\left|\dfrac{-5}{3}\right|>1$; $\left|\dfrac{1}{3}\right|<1$.\\ 
		 Vậy $\lim\left(\dfrac{1}{3}\right)^{n}=0$.
	 }
\end{ex}

% nguồn: Ok_Toan11_TracNghiem_Lan2
\begin{ex}%[1D3N1-4]%[ma:45d8f0dc]
	Tính giới hạn của  dãy số $u_n=\dfrac{-2+3n-2n^3}{3n-2}$.
	\choice 
		{$\dfrac{-2}{3}$} 
		 {\True $-\infty$}
		{$1$} 
		{$+\infty$} 
	\loigiai{ 
		Ta có $\lim\dfrac{-2+3n-2n^3}{3n-2}=\lim\left(n^2\cdot \dfrac{-\dfrac{2}{n^3}+\dfrac{3}{n^2}-2}{3-\dfrac{2}{n}}\right)=-\infty$ vì 
		 $\heva{&\lim n^2=+\infty \\&\lim\dfrac{-\dfrac{2}{n^3}+\dfrac{3}{n^2}-2}{3-\dfrac{2}{n}}=-\dfrac{2}{3}<0.}$
	 }
\end{ex}

% nguồn: Ok_Toan11_TracNghiem_Lan2
\begin{ex}%[1D3N1-4]%[ma:701313b2]
	Trong các giới  hạn sau, giới hạn nào bằng $0$?
	 \choice
		 {\True $\lim\left(\dfrac{2}{3}\right)^{n}$}
		 {$\lim\left(\dfrac{5}{3}\right)^{n}$}
		 {$\lim\left(\dfrac{4}{3}\right)^{n}$}
		 {$\lim 2^{n}$}
	 \loigiai{
		Áp dụng $\lim  q^{n}=0,\quad(|q|<1)$.
	} 
\end{ex}

% nguồn: Ok_Toan11_TracNghiem_Lan2
\begin{ex}%[1D3N1-4]%[ma:c2b7290a]
	Dãy số nào sau đây có giới  hạn bằng $0$?
	 \choice
		{\True  $(0{,}999)^{n}$}
		 {$(-1)^{n}$}
		 {$(-1{,}0001)^{n}$}
		{$(1{,}2345)^{n}$} 
	 \loigiai{
		Do $0{,}999<1$  nên $\lim (0{,}999)^{n}=0$.
	 }
\end{ex}

% nguồn: Ok_Toan11_TracNghiem_Lan2
\begin{ex}%[1D3N1-4]%[ma:a6db9252]
	Tính giới  hạn $L=\lim\left(\sqrt{4n^2+n+1}-9n\right)$.
	 \choice
		 {$+\infty$}
		{$-7$} 
		 {\True $-\infty$}
		{$\dfrac{9}{4}$} 
	 \loigiai{
		\begin{align*}
			L&=\lim\left(\sqrt{4  n^2+n+1}-9 n\right)\\
			 &=\lim\dfrac{\left(\sqrt{4 n^2+n+1}-9 n\right)\left(\sqrt{4 n^2+n+1}+9 n\right)}{\sqrt{4 n^2+n+1}+9 n}\\
			&=\lim\dfrac{4  n^2+n+1-81 n^2}{\sqrt{4 n^2+n+1}+9 n}\\
			&=\lim\dfrac{-77 n^2+n+1}{\sqrt{4  n^2+n+1}+9 n}\\
			&=\lim\dfrac{n^2\left(-77+\dfrac{1}{n}+\dfrac{1}{n^2}\right)}{n\left(\sqrt{4+\dfrac{1}{n}+\dfrac{1}{n^2}}+9\right)} \\ 
			& =\lim n \left(\dfrac{-77+\dfrac{1}{n}+\dfrac{1}{n^2}}{\sqrt{4+\dfrac{1}{n}+\dfrac{1}{n^2}}+9}\right)\\ 
			&=-\infty 
		\end{align*}
		Vì $\lim n=+\infty $ và  $\lim\left(\dfrac{-77+\dfrac{1}{n}+\dfrac{1}{n^2}}{\sqrt{4+\dfrac{1}{n}+\dfrac{1}{n^2}}+9} \right)=-7< 0$.
	} 
\end{ex}

% nguồn: Ok_Toan11_TracNghiem_Lan2
\begin{ex}%[1D3N1-4]%[ma:189f46a6]
	$\lim {\left(\dfrac{2\,018}{2\,019}\right)^{n}}$  bằng.
	\choice 
		{\True  $0$}
		 {$+\infty$}
		 {$\dfrac{1}{2}$}
		{$2$} 
	 \loigiai{
		Áp  dụng $\lim q^{n}=0,\quad|q|<1$.
	} 
\end{ex}

% nguồn: Ok_Toan11_TracNghiem_Lan2
\begin{ex}%[1D3N1-4]%[ma:b697c0f2]
	Tính giới hạn  $L=\lim\dfrac{n^3-2n}{3n^2+n-2}$.
	 \choice
		{\True $L=+\infty$} 
		 {$L=0$}
		 {$L=\dfrac{1}{3}$}
		{$L=-\infty$} 
	 \loigiai{
		Ta có $L=\lim\dfrac{n^3-2n}{3n^2+n-2}=\lim\left(n\cdot \dfrac{1-\dfrac{2}{n^2}}{3+\dfrac{1}{n}-\dfrac{2}{n^2}}\right)=+\infty  $ vì
		$\heva{&\lim  n=+\infty \\&\lim\dfrac{1-\dfrac{2}{n^2}}{3+\dfrac{1}{n}-\dfrac{2}{n^2}}=\dfrac{1}{3}>0.}$
	} 
\end{ex}

% nguồn: Ok_Toan11_TracNghiem_Lan2
\begin{ex}%[1D3N1-4]%[ma:69c09493]
	 Tính giới hạn $L=\lim\left(\sqrt{9n^2+2n-1}-\sqrt{4n^2+1} \right)$.
	\choice 
		 {\True $+\infty$}
		{$1$} 
		{$-\infty$} 
		 {$\dfrac{9}{4}$}
	\loigiai{ 
		\begin{align*}
			L&=\lim\left(\sqrt{9 n^2+2 n-1}-\sqrt{4  n^2+1}\right)\\
			&=\lim\dfrac{\left(\sqrt{9 n^2+2 n-1}-\sqrt{4  n^2+1}\right) \left(\sqrt{9 n^2+2 n-1}+\sqrt{4 n^2+1}\right)}{\sqrt{9 n^2+2 n-1}+\sqrt{4 n^2+1}}\\
			&=\lim\dfrac{\left(9 n^2+2 n-1\right)-\left(4 n^2+1\right)}{\sqrt{9 n^2+2 n-1}+\sqrt{4  n^2+1}}\\
			 &=\lim\dfrac{5 n^2+2 n-2}{\sqrt{9 n^2+2 n-1}+\sqrt{4 n^2+1}} \\
			 &=\lim\dfrac{n^2\left(5+\dfrac{2}{n}-\dfrac{2}{n^2}\right)}{n\left(\sqrt{9+\dfrac{2}{n}-\dfrac{1}{n^2}}+\sqrt{4+\dfrac{1}{n^2}}\right)}\\
			&=\lim  n \left(\dfrac{5+\dfrac{2}{n}-\dfrac{2}{n^2}}{\sqrt{9+\dfrac{2}{n}-\dfrac{1}{n^2}}+\sqrt{4+\dfrac{1}{n^2}}}\right)\\
			&=+\infty. 
		\end{align*}
	 }
\end{ex}

% nguồn: Ok_Toan11_TracNghiem_Lan2
\begin{ex}%[1D3N1-4]%[ma:725ed8d2]
	Giá trị đúng của $\lim\left(5^{n}\right)$  là
	 \choice
		{\True  $+\infty$}
		{$-2$} 
		{$2$} 
		 {$-\infty$}
	\loigiai{ 
		Ta  có $\lim\left(5^{n}\right)=\lim\dfrac{1}{\left(\dfrac{1}{5}\right)^{n}}=+\infty$ vì $\lim\left(\dfrac{1}{5}\right)^{n}=0$ và $\left(\dfrac{1}{5}\right)^{n}>0 \quad \forall n \in \mathbb{N}^{*}$.
	} 
\end{ex}

% nguồn: Ok_Toan11_TracNghiem_Lan2
\begin{ex}%[1D3N1-4]%[ma:25de2fda]
	$\lim _{n\rightarrow+\infty}2^{n}$  bằng
	 \choice
		 {$2$}
		{\True $+\infty$} 
		 {$-\infty$}
		 {$0$}
	 \loigiai{}
	Áp dụng $\lim q^{n}=+\infty,\quad(|q|>1)$. 
\end{ex}

% nguồn: Ok_Toan11_TracNghiem_Lan2
\begin{ex}%[1D3H1-4]%[ma:5f7137d9]
	Tìm  $\lim\dfrac{3n^4-2n+4}{4n^2+2n+3}$.
	 \choice
		{$-1$} 
		{\True $+\infty$} 
		 {$0$}
		{$\dfrac{3}{4}$} 
	\loigiai{Ta có $\lim\dfrac{3n^4-2n+4}{4n^2+2n+3}=\lim\dfrac{3-\dfrac{2}{n^3}+\dfrac{4}{n^4}}{\dfrac{4}{n^2}+\dfrac{2}{n^3}+\dfrac{3}{n^4}}=+\infty$. 
	 }
\end{ex}

% nguồn: Ok_Toan11_TracNghiem_Lan2
\begin{ex}%[1D3H1-4]%[ma:df5c7579]
	 Tính $\lim\left(2-3n\right)^4\left(n+1\right)^3$.
	 \choice
		 {$-\infty$}
		{\True  $+\infty$}
		{$81$} 
		{$2$} 
	\loigiai{ 
		Ta  có $\lim\left(2-3n\right)^4\left(n+1\right)^3=\lim\left[n^7\left(\dfrac{2}{n}-3\right)^4\left(1+\dfrac{1}{n} \right)^3\right]$.\\
		Lại  có $\lim n^7=+\infty$;
		$\lim\left(\dfrac{2}{n}-3\right)^4=\left(-3\right)^4=3^4$; 
		$\lim\left(1+\dfrac{1}{n} \right)^3=1$.\\ 
		Suy ra  $\lim\left(2-3n\right)^4\left(n+1\right)^3=+\infty$.
	} 
\end{ex}

% nguồn: Ok_Toan11_TracNghiem_Lan2
\begin{ex}%[1D3H1-4]%[ma:fb8fd153]
	Giới hạn  $\lim\left(n^2-n+1\right)(2-n)$ bằng
	\choice 
		 {$+\infty$}
		{\True  $-\infty$}
		{$0$} 
		 {$-1$}
	\loigiai{Ta có \begin{itemize}
		\item  $\lim\left(n^2-n+1\right)=\lim n^2\left(1-\dfrac{1}{n}+\dfrac{1}{n^2}\right)=+\infty$,
		\item $\lim (2-n)=\lim n\left(\dfrac{2}{n}-1\right)=-\infty$ 
	\end{itemize} nên $\lim\left(n^2-n+1\right)(-n+2)=-\infty$.
} 
\end{ex}

% nguồn: Ok_Toan11_TracNghiem_Lan2
\begin{ex}%[1D3H1-4]%[ma:6f887bec]
	Giới  hạn $\lim\left(2^{2n}-3^n\right)$ bằng
	 \choice
		 {$3$}
		 {$2$}
		 {$-\infty$}
		 {\True $+\infty$}
	 \loigiai{
	Ta có $\lim\left(2^{2n}-3^n\right)=\lim\left(4^n-3^n\right)=\lim 4^n\left[1-\left(\dfrac{3}{4}\right)^n\right]=+\infty$  vì $\lim 4^n=+\infty$ và $\lim\left[1-\left(\dfrac{3}{4}\right)^n\right]=1>0$.}
\end{ex}

% nguồn: Ok_Toan11_TracNghiem_Lan2
\begin{ex}%[1D3H1-4]%[ma:17e9f7af]
	Kết  quả của $\lim\dfrac{3^n-2\cdot 5^{2n}}{2^{n+1}+5^{n+5}}$ bằng
	\choice 
		{\True $-\infty$} 
		{$+\infty$} 
		{$2$} 
		{$\dfrac{3}{2}$} 
	\loigiai{ 
	 Ta có $\lim\dfrac{3^n-2.5^{2n}}{2^{n+1}+5^{n+5}}=\lim\dfrac{3^n-2.25^n}{2.2^n+5^5.5^n}=\lim\dfrac{\left(\dfrac{3}{5}\right)^n-2.5^n}{2.\left(\dfrac{2}{5}\right)^n+5^5}=-\infty$.}
\end{ex}

% nguồn: Ok_TracNghiemLop11_Lan1
\begin{ex}%[1D3H1-4]%[ma:c0c07193]
	Tính $\lim \dfrac{8n^2+n-2}{n^2}$. 
	\choice 
		 {$3$}
		 {$0$}
		 {$-2$}
		 {\True $8$}
	 \loigiai{
		Ta có  $\lim \dfrac{8n^2+n-2}{n^2}=\lim \left(8+\dfrac{1}{n}-\dfrac{2}{n^2}\right)=8$.
	} 
\end{ex}

% nguồn: Ok_TracNghiemLop11_Lan1
\begin{ex}%[1D3H1-4]%[ma:40b8c04a]
	Kết quả của giới hạn  $\lim \sqrt{{{2\cdot 3}^n}-n+2}$ là
	\choice 
		 {$0$}
		{$2$} 
		 {$3$}
		{\True $+\infty$} 
	 \loigiai{
		Ta có $\lim \sqrt{{{2\cdot  3}^n}-n+2}=\lim \sqrt{{3^n}}\cdot\sqrt{2-\dfrac{n}{{3^n}}+2\cdot{{\left(\dfrac{1}{3} \right)}^n}}$.\\
		 Vì
		$\left\{ \begin{aligned}
		& \lim \sqrt{{3^n}}=+\infty \\ & 0\le \dfrac{n}{{3^n}}\le \dfrac{n}{C_n^2}=\dfrac{2}{n-1}\to 0\Rightarrow \lim \dfrac{n}{{3^n}}=0\\ & \lim {{\left(\dfrac{1}{3} \right)}^n}=0\\ 
	\end{aligned} \right\}\xrightarrow[{}]{}\left\{\begin{aligned}
	& \lim \sqrt{{3^n}}=+\infty \\  & \lim \sqrt{2-\dfrac{n}{{3^n}}+2\cdot {{\left(\dfrac{1}{3} \right)}^n}}=\sqrt{2} > 0\\
	\end{aligned} \right.$,
	do đó $\lim  \sqrt{{{2\cdot 3}^n}-n+2} = +\infty$.
	} 
\end{ex}

% nguồn: Ok_Toan11_TracNghiem_Lan2
\begin{ex}%[1D3V1-4]%[ma:9cfe1da6]
	Chọn kết  quả đúng của $\lim\dfrac{\sqrt{n^2-2n+5}-\sqrt{n^4+3n^2+1}}{3+5n}$.
	\choice 
		{\True  $-\infty$}
		 {$0$}
		{$\dfrac{1}{5}$} 
		 {$+\infty$}
	\loigiai{ 
		Ta có  $\lim\dfrac{\sqrt{n^2-2n+5}-\sqrt{n^4+3n^2+1}}{3+5n}=\lim\dfrac{\sqrt{\dfrac{1}{n^2}-\dfrac{1}{n^3}+\dfrac{5}{n^4}}-\sqrt{1+\dfrac{3}{n^2}+\dfrac{1}{n^4}}}{\dfrac{3}{n^2}+\dfrac{5}{n}}$.\\
		Từ \begin{itemize}
		\item  $\lim\left(\sqrt{\dfrac{1}{n^2}-\dfrac{1}{n^3}+\dfrac{5}{n^4}}-\sqrt{1+\dfrac{3}{n^2}+\dfrac{1}{n^4}}\right)=-1$,
		 \item $\lim\left(\dfrac{3}{n^2}+\dfrac{5}{n}\right)=0$ và $\dfrac{3}{n^2}+\dfrac{5}{n}>0$, $\forall n>0$
	\end{itemize}
	suy ra $\lim\dfrac{\sqrt{\dfrac{1}{n^2}-\dfrac{1}{n^3}+\dfrac{5}{n^4}}-\sqrt{1+\dfrac{3}{n^2}+\dfrac{1}{n^4}}}{\dfrac{3}{n^2}+\dfrac{5}{n}}=-\infty$. 
} 
\end{ex}

% nguồn: Ok_Toan11_TracNghiem_Lan2
\begin{ex}%[1D3V1-4]%[ma:448e9b47]
	Kết quả  của $\lim n\left(\sqrt{n+5}-\sqrt{n}\right)$ bằng
	 \choice
		{\True $+\infty$} 
		 {$-\infty$}
		{$2$} 
		{$0$} 
	 \loigiai{
		Ta có $\lim  n\left(\sqrt{n+5}-\sqrt{n}\right)=\lim\dfrac{5n}{\sqrt{n+5}+\sqrt{n}}=\lim\dfrac{5}{\sqrt{\dfrac{1}{n^2}+\dfrac{5}{n}}+\sqrt{\dfrac{1}{n}}}$.\\
		Vì $\lim 5=5>0$, $\lim\left(\sqrt{\dfrac{1}{n}+\dfrac{5}{n^2}}+\sqrt{\dfrac{1}{n}}\right)=0$ và $\sqrt{\dfrac{1}{n}+\dfrac{5}{n^2}}+\sqrt{\dfrac{1}{n}}>0$,  $\forall n\in\mathbb{N}^*$ nên $$\lim n\left(\sqrt{n+5}-\sqrt{n}\right)=+\infty.$$
	 }
\end{ex}

% nguồn: Ok_Toan11_TracNghiem_Lan2
\begin{ex}%[1D3V1-4]%[ma:5e0f50d2]
	 $\lim\left[\sqrt{2n^2+2n+1}-\sqrt[3]{n^3+n^2}\right]$ bằng
	\choice 
		{$-\infty$} 
		{$0$} 
		 {\True $+\infty$}
		 {$2$}
	 \loigiai{
		Ta có $\lim\left[\sqrt{2n^2+2n+1}-\sqrt[3]{n^3+n^2}\right]=\lim  n\left[\sqrt{2+\dfrac{2}{n}+\dfrac{1}{n^2}}-\sqrt[3]{1+\dfrac{1}{n}}\right]=+\infty$
		vì $\lim n=+\infty$  và $\lim\left[\sqrt{2+\dfrac{2}{n}+\dfrac{1}{n^2}}-\sqrt[3]{1+\dfrac{1}{n}}\right]=\sqrt{2}-1>0$.
	} 
\end{ex}

% nguồn: Ok_Toan11_TracNghiem_Lan2
\begin{ex}%[1D3V1-4]%[ma:c7b2ca8a]
	Giá  trị của $E=\lim\left(\sqrt{n^2+n+1}-\sqrt[3]{2n^3-3n^2+1}\right)$ bằng
	 \choice
		{$+\infty$} 
		{\True  $-\infty$}
		{$0$} 
		 {$1$}
	\loigiai{ Ta  có $E=\lim\left(\sqrt{n^2+n+1}-\sqrt[3]{2n^3-3n^2+1}\right)=\lim n\left(\sqrt{1+\dfrac{1}{n}+\dfrac{1}{n^2}}-\sqrt[3]{2-\dfrac{3}{n}+\dfrac{1}{n^3}}\right)$.\\
		Vì $\lim n=+\infty$ và $\lim\left(\sqrt{1+\dfrac{1}{n}+\dfrac{1}{n^2}}-\sqrt[3]{2-\dfrac{3}{n}+\dfrac{1}{n^3}}\right)=1-\sqrt[3]{2}<0$ nên  $$E=\lim\left(\sqrt{n^2+n+1}-\sqrt[3]{2n^3-3n^2+1}\right)=-\infty.$$
	 }
\end{ex}


%%%%%%%%%%%%%%%% ĐÚNG/SAI %%%%%%%%%%%%%%%%

% nguồn: Ok_dungsai_lop11_Lan1
\begin{ex}%[1D3N1-4]%[ma:9071321e]
 Cho các hàm số $f(x)$, $g(x)$, $h(x)$ thỏa mãn $\lim\limits_{x \to x_0} f(x)=2$, $\lim\limits_{x \to x_0} g(x)=5$, $\lim\limits_{x \to  x_0} h(x)=0$. Mỗi kết quả giới hạn dưới đây đúng hay sai?
  \choiceTF
  {$\lim\limits_{x  \to x_0} (g(x)-f(x))=-3$}
  {\True $\lim\limits_{x \to  x_0} \dfrac{f(x)}{g(x)}=\dfrac{2}{5}$}
   {\True $\lim\limits_{x \to x_0} \dfrac{h(x)}{g(x)}=0$}
  {\True $\lim\limits_{x \to  x_0} (g(x)\cdot h(x))=0$}
  \loigiai{
  \begin{itemchoice}
   \itemch $\lim\limits_{x \to x_0} (g(x)-f(x))=5-2=3$. 
   \itemch  $\lim\limits_{x \to x_0} \dfrac{f(x)}{g(x)}=\dfrac{2}{5}$.
   \itemch  $\lim\limits_{x \to x_0} \dfrac{h(x)}{g(x)}=\dfrac{0}{5}=0$.
   \itemch $\lim\limits_{x \to x_0} (g(x)\cdot h(x))=5\cdot 0=0$. 
  \end{itemchoice}
  }
\end{ex}

% nguồn: Ok_dungsai_lop11_Lan1
\begin{ex}%[1D3N1-4]%[ma:5c108d03]
 Mỗi kết quả giới hạn dưới  đây đúng hay sai?
  \choiceTF
  {\True $\lim\limits_{n \to  +\infty} \left(n+1\right)=+\infty$}
   {\True $\lim\limits_{n \to +\infty} (-n^2-n+2)=-\infty$}
  {\True $\lim\limits_{n \to +\infty}  \dfrac{2n^2+2n-1}{n+3}=+\infty$}
  {$\lim\limits_{n \to +\infty} \dfrac{2^n+3^n}{2^n+1}=1$} 
 \loigiai{ 
  \begin{itemchoice}
   \itemch $\lim\limits_{n \to +\infty}  \left(n+1\right)
   =\lim\limits_{n \to +\infty} n\left(1+\dfrac{1}{n}\right)=+\infty$, vì $\lim\limits_{n \to +\infty} n=+\infty$, $\lim\limits_{n \to  +\infty} \left(1+\dfrac{1}{n}\right)=1>0$.
    \itemch $\lim\limits_{n \to +\infty} (-n^2-n+2)
    =\lim\limits_{n \to +\infty} n^2\left(-1-\dfrac{1}{n}+\dfrac{2}{n^2}\right)=-\infty
    $,\\
   vì $\lim\limits_{n \to  +\infty} n^2=+\infty$, $\lim\limits_{n \to +\infty} \left(-1-\dfrac{1}{n}+\dfrac{2}{n^2}\right)=-1<0$.
   \itemch $\lim\limits_{n \to  +\infty} \dfrac{2n^2+2n-1}{n+3}=\lim\limits_{n \to +\infty} \dfrac{n^2\left(2+\dfrac{2}{n}-\dfrac{1}{n^2}\right)}{n\left(1+\dfrac{3}{n}\right)}
    =\lim\limits_{n \to +\infty} n\cdot \dfrac{2+\dfrac{2}{n}-\dfrac{1}{n^2}}{1+\dfrac{3}{n}}
    =+\infty$,\\
    vì $\lim\limits_{n \to +\infty} n=+\infty$, $\lim\limits_{n \to +\infty} \dfrac{2+\dfrac{2}{n}-\dfrac{1}{n^2}}{1+\dfrac{3}{n}}=2>0$.
   \itemch  $\lim\limits_{n \to +\infty} \dfrac{2^n+3^n}{2^n+1}
   =\lim\limits_{n \to +\infty}  \dfrac{3^n\cdot \left(\left(\dfrac{2}{3}\right)^n+1\right)}{2^n\left(1+\left(\dfrac{1}{2}\right)^n\right)}
   =\lim\limits_{n \to  +\infty} \left(\dfrac{3}{2}\right)^n\cdot \dfrac{\left(\dfrac{2}{3}\right)^n+1}{1+\left(\dfrac{1}{2}\right)^n}
    =+\infty$,\\
   vì $\lim\limits_{n \to +\infty} \left(\dfrac{3}{2}\right)^n=+\infty$,  $\lim\limits_{n \to +\infty} \dfrac{\left(\dfrac{2}{3}\right)^n+1}{1+\left(\dfrac{1}{2}\right)^n}=1>0$.
  \end{itemchoice}
  }
\end{ex}

% nguồn: Ok_dungsai_lop11_Lan1
\begin{ex}%[1D3H1-4]%[ma:2b60bcec]
 Mỗi mệnh đề  sau đây đúng hay sai?
  \choiceTF
  {Cấp số nhân lùi vô hạng có số hạng đầu $u_1=1$, công bội $q=-\dfrac{1}{2}$ có tổng $S=2$} 
  {Cấp số nhân lùi vô hạn có tổng $S=4$, công bội  $q=\dfrac{1}{2}$ có số hạng đầu $u_1=1$}
  {\True Cấp số nhân lùi vô hạn có tổng  $S=-6$, số hạng đầu $u_1=-3$ có công bội $q=\dfrac{1}{2}$}
  {\True Số thập phân vô hạn tuần hoàn $1{,}0202(02)$ viết dưới dạng  phân thức là $\dfrac{101}{99}$}
 \loigiai{ 
  Cấp số nhân lùi vô hạn có số hạng đầu  $u_1$, công bội $q$ có tổng $S=\dfrac{u_1}{1-q}$. Suy ra $u_1=(1-q)\cdot S$, $q=1-\dfrac{u_1}{S}$.
  \begin{itemchoice}
   \itemch Cấp số nhân lùi vô hạng  có số hạng đầu $u_1=1$, công bội $q=-\dfrac{1}{2}$ có tổng $S=\dfrac{1}{1+\dfrac{1}{2}}=\dfrac{2}{3}$.
   \itemch Cấp  số nhân lùi vô hạn có tổng $S=4$, công bội $q=\dfrac{1}{2}$ có số hạng đầu $u_1=\left(1-\dfrac{1}{2}\right)\cdot 4=2$.
   \itemch Cấp số nhân lùi vô hạn có tổng $S=-6$, số hạng đầu $u_1=-3$ có công bội $q=1-\dfrac{-3}{-6}=\dfrac{1}{2}$. 
   \itemch Ta có  $1{,}0202(02)=1+\dfrac{2}{10^2}+\dfrac{2}{10^4}+\cdots+\dfrac{2}{10^{2n}}+\cdots
   =1+\dfrac{\dfrac{2}{10^2}}{1-\dfrac{1}{10^2}}=\dfrac{101}{99}$. 
  \end{itemchoice}
  }
\end{ex}

% nguồn: Ok_dungsai_lop11_Lan1
\begin{ex}%[1D3H1-4]%[ma:85e093bb]
 Mỗi kết quả giới hạn dưới đây đúng hay  sai?
 \choiceTF 
   {$\lim \limits_{n\rightarrow +\infty} \dfrac{3^{n+2}+2^n}{3^{n+1} +1}=+\infty$}
   {$\lim\limits_{n \to +\infty} \sqrt[3]{2+n-n^3}=+\infty$}
   {$\lim\limits_{n \to +\infty} \dfrac{-n^3+n^2+2}{n+5}=+\infty$}
  {\True $\lim\limits_{n \to  +\infty} \left(\sqrt{n^2+2}+n\right)=+\infty$}
 \loigiai{ 
  \begin{itemchoice}
    \itemch $\lim \limits_{n\rightarrow +\infty} \dfrac{3^{n+2}+2^n}{3^{n+1} +1} =\lim \limits_{n\rightarrow +\infty} \dfrac{9 \cdot 3^n+2^n}{3\cdot 3^n+1} =\lim \limits_{n\rightarrow +\infty} \dfrac{9+\left(\dfrac{2}{3}\right)^n}{3+\left(\dfrac{1}{3}\right)^n} = 3$.
   \itemch $\lim\limits_{n \to +\infty} \sqrt[3]{2+n-n^3} 
   =\lim\limits_{n \to  +\infty} n\sqrt[3]{\dfrac{2}{n^3}+\dfrac{1}{n^2}-1}=-\infty
   $,\\ 
   vì $\lim\limits_{n \to +\infty} n=+\infty$, $\lim\limits_{n \to +\infty}  \sqrt[3]{\dfrac{2}{n^3}+\dfrac{1}{n^2}-1}=-1<0$.
   \itemch $\lim\limits_{n \to +\infty} \dfrac{-n^3+n^2+2}{n+5} 
   =\lim\limits_{n \to +\infty} \dfrac{n^3\left(-1+\dfrac{1}{n}+\dfrac{2}{n^3}\right)}{n\left(1+\dfrac{5}{n}\right)} 
   =\lim\limits_{n \to +\infty}  n^2\cdot \dfrac{-1+\dfrac{1}{n}+\dfrac{2}{n^3}}{1+\dfrac{5}{n}}
    =-\infty$,\\
   vì $\lim\limits_{n  \to +\infty} n^2=+\infty$, $\lim\limits_{n \to +\infty} \dfrac{-1+\dfrac{1}{n}+\dfrac{2}{n^3}}{1+\dfrac{5}{n}}=-1<0$.
    \itemch $\lim\limits_{n \to +\infty} \left(\sqrt{n^2+2}+n\right)
   =\lim\limits_{n  \to +\infty} n\left(\sqrt{1+\dfrac{2}{n^2}}+1\right)=+\infty
   $,\\ 
   vì $\lim\limits_{n \to +\infty} n=+\infty$, $\lim\limits_{n \to +\infty} \left(\sqrt{1+\dfrac{2}{n^2}}+1\right)=2>0$. 
  \end{itemchoice}
  }
\end{ex}

% nguồn: Ok_dungsai_lop11_Lan1
\begin{ex}%[1D3H1-4]%[ma:5d326372]
 Xét tính đúng sai của các mệnh đề sau: 
  \choiceTF
  {$\lim  (\sqrt{3})^n=-\infty$}
  {$\lim  \pi^n=0$}
   {\True $\lim \left(n^3+2n^2-4\right)=+\infty$}
  {\True $\lim \left(-n^4+5n^3-4n\right)=-\infty$} 
  \loigiai{
  \begin{itemchoice}
   \itemch Sai. $\lim (\sqrt{3})^n=+\infty ($ do $\sqrt{3} >  1)$.
   \itemch Sai. $\lim  \pi^n=+\infty ($ do $\pi > 1)$.
   \itemch Đúng. $\lim \left(n^3+2n^2-4\right)=\lim n^3 \cdot  \left(1+\dfrac{2}{n}-\dfrac{4}{n^3} \right)$.\\
   Vì  $\heva{&{\lim n^3=+\infty} \\
   &{\lim \left(1+\dfrac{2}{n}-\dfrac{4}{n^3} \right)=1>  0}
   }$  nên $\lim \left(n^3+2n^2-4\right)=+\infty$.
   \itemch Đúng. $\lim \left(-n^4+5n^3-4n\right)=\lim n^4 \cdot  \left(-1+\dfrac{5}{n}-\dfrac{4}{n^3} \right)$.\\
    Vì $\heva{&{\lim n^4=+\infty} \\
   &{\lim  \left(-1+\dfrac{5}{n}-\dfrac{4}{n^3} \right)=-1< 0}
   }$ nên $\lim  \left(-n^4+5n^3-4n\right)=-\infty$.
  \end{itemchoice}
  }
\end{ex}

% nguồn: Ok_dungsai_lop11_Lan2
\begin{ex}%[1D3H1-4]%[ma:beb7d4fa]
	Tính được các giới hạn. Xét tính đúng sai của  các khẳng định sau:
	 \choiceTF
		 {$\lim \left(\sqrt{2}\right)^n=-\infty$}
		{$\lim \pi^n=0$} 
		{\True $\lim  \left(2n^3+2n^2-4\right)=+\infty$}
		{\True $\lim  \left(-5n^4+n^3-4n\right)=-\infty$}
	 \loigiai{
		\begin{itemchoice}
			 \itemch Ta có $\lim \left(\sqrt{2}\right)^n=+\infty$ (do $\sqrt{2} > 1$).
			\itemch Ta có $\lim \pi^n=+\infty$ (do $\pi > 1$). 
			\itemch  Ta có $\lim \left(2n^3+2n^2-4\right)=\lim \left[n^3\cdot\left(2+\dfrac{2}{n}-\dfrac{4}{n^3}\right)\right]=+\infty$.\\
			Vì $\heva{&\lim n^3=+\infty  \\&\lim \left(2+\dfrac{2}{n}-\dfrac{4}{n^3}\right)=2> 0.}$
			\itemch Vì  $\lim \left(-5n^4+n^3-4n\right)=\lim \left[n^4\cdot\left(-5+\dfrac{1}{n}-\dfrac{4}{n^3}\right)\right]=-\infty$.\\
			 Vì $\heva{&\lim n^4=+\infty \\&\lim \left(-5+\dfrac{1}{n}-\dfrac{4}{n^3}\right)=-5< 0.}$
		\end{itemchoice}
	} 
\end{ex}


%%%%%%%%%%%%%%%% TỰ LUẬN %%%%%%%%%%%%%%%%

% nguồn: Ok_TuanLuan_Toan11_Lan2
\begin{ex}%[1D3H1-4]%[ma:a77542a8]
	Tính $\lim \dfrac{(2  n-1)\left(3 n^{2}+2\right)^{3}}{-2 n^{5}+4 n^{3}-1}$.
	\loigiai{Ta  có
		$\lim \dfrac{(2 n-1)\left(3 n^{2}+2\right)^{3}}{-2 n^{5}+4 n^{3}-1}=\lim  \dfrac{n^{7}\left(2-\dfrac{1}{n}\right)\left(3+\dfrac{2}{n^{2}}\right)^{3}}{n^{5}\left(-2+\dfrac{4}{n^{2}}-\dfrac{1}{n^{5}}\right)}=\lim n^{2} \dfrac{\left(2-\dfrac{1}{n}\right)\left(3+\dfrac{2}{n^{2}}\right)^{3}}{-2+\dfrac{4}{n^{2}}-\dfrac{1}{n^{5}}}=-\infty$\\
		(vì $\lim n^{2}=+\infty $; $\lim n^{2}  \dfrac{\left(2-\dfrac{1}{n}\right)\left(3+\dfrac{2}{n^{2}}\right)^{3}}{-2+\dfrac{4}{n^{2}}-\dfrac{1}{n^{5}}}=-27<0$).
	 }
\end{ex}

% nguồn: Ok_TuanLuan_Toan11_Lan2
\begin{ex}%[1D3H1-4]%[ma:23a65254]
	 Tính $\lim \left(n^{4}-2 n^{2}+3\right)$.
	\loigiai{Ta có 
		$\lim \left(-2 n^{3}+3 n-1\right)=\lim \left[n^{3}\left(-2+\dfrac{3}{n^{2}}-\dfrac{1}{n^{3}}\right)\right]  =-\infty$\\
		 (vì $\lim n^{3}=+\infty $; $\lim \left(-2+\dfrac{3}{n^{2}}-\dfrac{1}{n^{3}}\right)=-2<0$).
	 }
\end{ex}

% nguồn: Ok_TuanLuan_Toan11_Lan2
\begin{ex}%[1D3H1-4]%[ma:5368e97f]
	Tính $\lim \left(-2 n^{2}+4\right)^{3}$. 
	 \loigiai{Ta có
		$\lim \left(-2  n^{2}+4\right)^{3}=\lim n^{6}\left(-2+\dfrac{4}{n^{2}}\right)^{3}=-\infty$\\
		(vì $\lim n^{6}=+\infty $; $ \lim \left(-2+\dfrac{4}{n^{2}}\right)^{3}=-8<0$). 
	} 
\end{ex}

% nguồn: Ok_TuanLuan_Toan11_Lan2
\begin{ex}%[1D3H1-4]%[ma:20bb987e]
	Tính  $\lim \left(2 n-\sqrt{n^{3}+2 n-2}\right)$.
	\loigiai{Ta có $\lim \left(2 n-\sqrt{n^{3}+2 n-2}\right)=\lim\left[n  \sqrt{n}\left(\dfrac{2}{\sqrt{n}}-\sqrt{1+\dfrac{2}{n^{2}}-\dfrac{2}{n^{3}}}\right)\right] =-\infty$\\
		(vì $\lim  n \sqrt{n}=+\infty $; $\lim \left(\dfrac{2}{\sqrt{n}}-\sqrt{1+\dfrac{2}{n^{2}}-\dfrac{2}{n^{3}}}\right)=-1<0$).
	} 
\end{ex}

% nguồn: Ok_TuanLuan_Toan11_Lan2
\begin{ex}%[1D3H1-4]%[ma:4bb6be2d]
	Tính $\lim \dfrac{3 n^{2}-\sqrt{2 n^{4}+3 n-2}}{4  n-\sqrt{3 n^{2}+2}}$.
	\loigiai{Ta  có
		\begin{eqnarray*}
			\lim \dfrac{3 n^{2}-\sqrt{2 n^{4}+3 n-2}}{4 n-\sqrt{3 n^{2}+2}}&=&\lim \dfrac{n^{2}\left(3-\sqrt{2+\dfrac{3}{n^{3}}-\dfrac{2}{n^{4}}}\right)}{n\left(4-\sqrt{3+\dfrac{2}{n^{2}}}\right)}\\ 
			&=&\lim \left[n  \dfrac{\left(3-\sqrt{2+\dfrac{3}{n^{3}}-\dfrac{2}{n^{4}}}\right)}{\left(4-\sqrt{3+\dfrac{2}{n^{2}}}\right)}\right]\\
			 &=&+\infty
		\end{eqnarray*}
		 (vì $\lim n=+\infty $; $\lim \dfrac{\left(3-\sqrt{2+\dfrac{3}{n^{3}}-\dfrac{2}{n^{4}}}\right)}{\left(4-\sqrt{3+\dfrac{2}{n^{2}}}\right)}=\dfrac{3-\sqrt{2}}{4-\sqrt{3}}>0$).
	} 
\end{ex}

% nguồn: Ok_TuanLuan_Toan11_Lan2
\begin{ex}%[1D3H1-4]%[ma:af90e4cf]
	Tính $\lim  \left(5^{n}-2^{n}\right)$.
	 \loigiai{Ta có $\lim \left(5^{n}-2^{n} \right)=\lim \left[5^{n}\left(1-\left(\dfrac{2}{5}\right)^{n}\right)\right]=+\infty$\\
		(vì $\lim 5^{n}=+\infty$ và $\lim  \left(1-\left(\dfrac{2}{5}\right)^{n}\right)=1>0$).
	} 
\end{ex}

% nguồn: Ok_TuanLuan_Toan11_Lan2
\begin{ex}%[1D3H1-4]%[ma:23f04170]
	Cho dãy số $\left(u_n\right)$ xác định $u_{1}=0$, $u_{2}=1$, $u_{n+1}=2$, $u_n-u_{n-1}+2$  với mọi $n \geq 2$. Tìm giới hạn của dãy số $\left(u_n\right)$.
	\loigiai{Ta có $u_{1}=0, u_{2}=1, u_{3}=4, u_{4}=9$.\\ 
		Bằng phương pháp quy nạp, dễ dàng chứng minh được $u_n=(n-1)^{2}$ với mọi $n  \geq 1$.\\ Do đó $\lim u_n=\lim (n-1)^{2}=+\infty$.
	 }
\end{ex}

% nguồn: Ok_TuanLuan_Toan11_Lan2
\begin{ex}%[1D3H1-4]%[ma:8708a24f]
	Tính $\lim \left(3.2^{n+1}-5.3^{n}+7  n\right)$.
	 \loigiai{Ta có
		 $\lim \left(3.2^{n+1}-5.3^{n}+7 n\right)=\lim \left[3^{n}\left(-5+6\left(\dfrac{2}{3}\right)^{n}+7 \dfrac{n}{3^{n}}\right) \right]=-\infty$\\
		(vì  $\lim 3^n=+\infty$; $\lim \left(-5+6\left(\dfrac{2}{3}\right)^{n}+7 \dfrac{n}{3^{n}}\right)=-5<0$).
	} 
\end{ex}

% nguồn: Ok_TuanLuan_Toan11_Lan2
\begin{ex}%[1D3H1-4]%[ma:eb31e47a]
	Tính $\lim  \dfrac{9^{n}-3.4^{n}}{6.7^{n}+8^{n}}$.
	\loigiai{Ta  có
		$\lim \dfrac{9^{n}-3 \cdot 4^{n}}{6 \cdot 7^{n}+8^{n}}=\lim \dfrac{9^{n}\left[1-3 \cdot\left(\dfrac{4}{9}\right)^{n}\right]}{8^{n}\left[6\left(\dfrac{7}{8}\right)^{n}+1\right]}=\lim \left[\left(\dfrac{9}{8}\right)^{n} \dfrac{\left(1-3  \cdot\left(\dfrac{4}{9}\right)^{n}\right)}{\left(6\left(\dfrac{7}{8}\right)^{n}+1\right)}\right]=+\infty$\\
		 (vì $\lim \left(\dfrac{9}{8}\right)^{n}=+\infty$; $\lim \left[\dfrac{1-3 \cdot\left(\dfrac{4}{9}\right)^{n}}{6\left(\dfrac{7}{8}\right)^{n}+1}\right]=1$).
	} 
\end{ex}

% nguồn: Ok_TuanLuan_Toan11_Lan2
\begin{ex}%[1D3H1-4]%[ma:0e98af3e]
	Tính $\lim \dfrac{2  n^{4}-3 n^{3}+2}{n^{3}+2}$.
	\loigiai{Ta có 
		$\lim \dfrac{2 n^{4}-3 n^{3}+2}{n^{3}+2}=\lim \dfrac{n^{4}\left(2-\dfrac{3}{n^{2}}+\dfrac{2}{n^{4}}\right)}{n^{3}\left(1+\dfrac{2}{n^{3}}\right)}=\lim \left[n  \cdot \dfrac{2-\dfrac{3}{n^{2}}+\dfrac{2}{n^{4}}}{1+\dfrac{2}{n^{3}}}\right]=+\infty$\\
		(vì $\lim n=+\infty $; $\lim  \dfrac{2-\dfrac{3}{n^{2}}+\dfrac{2}{n^{4}}}{1+\dfrac{2}{n^{3}}}=2>0$).
	} 
\end{ex}

% nguồn: Ok_TuanLuan_Toan11_Lan2
\begin{ex}%[1D3H1-4]%[ma:7255c96b]
	 Tính $\lim \dfrac{3+3^{2}+3^{3}+\ldots+3^{n}}{1+2+2^{2}+\ldots+2^{n}}$.
	\loigiai{Ta  có tử thức là tổng của $n$ số hạng đầu tiên của cấp số nhân $\left(u_n\right)$ với $u_{1}=3$ và công bội $q=3$.\\
		 Do đó $3+3^{2}+\ldots+3^{n}=\dfrac{3\left(3^{n}-1\right)}{2}$.\\
		 Ta có mẫu thức là tổng của $n+1$ số hạng đầu tiên của cấp số nhân $\left(v_n\right)$ với $v_{1}=1$ và công bội $q=2$.\\
		Do đó $1+2+2^{2}+\ldots+2^{n}=\dfrac{2\left(2^{n+1}-1\right)}{1}$.\\ 
		Vậy $\lim  \dfrac{3+3^{2}+3^{3}+\ldots+3^{n}}{1+2+2^{2}+\ldots+2^{n}}=\lim \left[\left(\dfrac{3}{2}\right)^{n} \cdot \dfrac{3\left(1-\left(\dfrac{1}{3}\right)^{n}\right)}{4\left(2-\left(\dfrac{1}{2}\right)^{n}\right)}\right]=+\infty$\\
		(vì $\lim \left(  \dfrac{3}{2}\right)^n = +\infty$; $\lim \dfrac{3\left(1-\left(\dfrac{1}{3}\right)^{n}\right)}{4\left(2-\left(\dfrac{1}{2}\right)^{n}\right)}=\dfrac{3}{8}>0$).
	} 
\end{ex}

% nguồn: Ok_TuanLuan_Toan11_Lan2
\begin{ex}%[1D3H1-4]%[ma:dc758247]
	 Tính $\lim \left(n^{2}-n \sqrt{4 n+1}\right)$.
	\loigiai{Ta  có $\lim \left(n^{2}-n \sqrt{4 n+1} \right)=\lim\left[n^{2}\left(1-\sqrt{\dfrac{4}{n}+\dfrac{1}{n^{2}}}\right)\right]=+\infty$\\
	(vì $\lim n^{2}=+\infty$;  $\lim \left(1-\sqrt{\dfrac{4}{n}+\dfrac{1}{n^{2}}}\right)=1>0$).}
\end{ex}
